% Derivazioni migrate dal Cap. 5 per il vincolo di lunghezza del corpo.
% Il capitolo mantiene il canale depolarizzante in forma completa e gli enunciati
% compatti degli altri tre.

I quattro canali quantistici che modellano il rumore sono mappe \ac{CPTP} in rappresentazione di Kraus, secondo l'Eq.~\ref{eq:kraus}. Il Capitolo~\ref{chap:rumore} ne riporta il canale depolarizzante per esteso, essendo quello su cui poggia la campagna sperimentale, e l'enunciato compatto degli altri tre; qui se ne danno le derivazioni complete.

\subsubsection{Canale Bit Flip}

Il canale \textbf{bit flip} modella un errore di tipo $X$ discreto: con probabilità $p$ il qubit viene invertito ($\ket{0}\leftrightarrow\ket{1}$), con probabilità $1-p$ rimane invariato. Gli operatori di Kraus sono $K_0 = \sqrt{1-p}\,I$ e $K_1 = \sqrt{p}\,X$, e il canale agisce come:
\begin{equation}
    \mathcal{E}_{BF}(\rho) = (1-p)\,\rho + p\,X\rho X
    \label{eq:bitflip}
\end{equation}
Sulla sfera di Bloch, le componenti $r_y$ e $r_z$ vengono contratte di un fattore $(1-2p)$ mentre la componente $r_x$ rimane invariata. Per $p=1/2$ lo stato è completamente distrutto sul piano $yz$.

% ---
\subsubsection{Canale Phase Flip}

Il canale \textbf{phase flip} modella la perdita di informazione di fase senza cambiamento di popolazione: la sovrapposizione decade senza che l'energia venga dissipata. Con operatori di Kraus $K_0 = \sqrt{1-p}\,I$ e $K_1 = \sqrt{p}\,Z$, il canale è:
\begin{equation}
    \mathcal{E}_{PF}(\rho) = (1-p)\,\rho + p\,Z\rho Z
    \label{eq:phaseflip}
\end{equation}
L'effetto sulla matrice densità è immediato:
\begin{equation}
    \mathcal{E}_{PF}\!\begin{pmatrix}\rho_{00} & \rho_{01} \\ \rho_{10} & \rho_{11}\end{pmatrix}
    = \begin{pmatrix}\rho_{00} & (1-2p)\,\rho_{01} \\ (1-2p)\,\rho_{10} & \rho_{11}\end{pmatrix}
\end{equation}
Le popolazioni $\rho_{00}$ e $\rho_{11}$ rimangono invariate, mentre le \textbf{coerenze} $\rho_{01}$, $\rho_{10}$ vengono soppresse di un fattore $(1-2p)$. Questo è il canale che descrive il decadimento $T_2$: parametrizzando $p(t) = \frac{1}{2}(1-e^{-t/T_2})$, si ottiene $\rho_{01}(t) = \rho_{01}(0)\,e^{-t/T_2}$.

% ---
\subsubsection{Canale di Amplitude Damping}

Il canale di \textbf{amplitude damping} modella il rilassamento energetico $T_1$: il qubit decade spontaneamente dallo stato eccitato $\ket{1}$ allo stato fondamentale $\ket{0}$ con probabilità $\gamma = 1-e^{-t/T_1}$. Gli operatori di Kraus sono:
\begin{equation}
    K_0 = \begin{pmatrix}1 & 0 \\ 0 & \sqrt{1-\gamma}\end{pmatrix}, \qquad
    K_1 = \begin{pmatrix}0 & \sqrt{\gamma} \\ 0 & 0\end{pmatrix}
    \label{eq:kraus_ad}
\end{equation}
Il canale agisce sulla matrice densità come:
\begin{equation}
    \mathcal{E}_{AD}\!(\rho) = \begin{pmatrix}
        \rho_{00} + \gamma\,\rho_{11} & \sqrt{1-\gamma}\,\rho_{01} \\
        \sqrt{1-\gamma}\,\rho_{10}   & (1-\gamma)\,\rho_{11}
    \end{pmatrix}
    \label{eq:amp_damping_explicit}
\end{equation}
La popolazione $\rho_{11}$ decade verso $\rho_{00}$ con tasso $\gamma$, mentre le coerenze vengono soppresse di $\sqrt{1-\gamma}$. A differenza dei canali precedenti, questo canale è \textbf{non simmetrico}: sulla sfera di Bloch contrae la sfera verso il polo nord $\ket{0}$, non verso il centro:
\begin{equation}
    r_x \to \sqrt{1-\gamma}\,r_x, \qquad r_y \to \sqrt{1-\gamma}\,r_y, \qquad r_z \to (1-\gamma)\,r_z + \gamma
\end{equation}
Per $t \to \infty$ ($\gamma \to 1$), qualsiasi stato converge a $\ket{0}\bra{0}$, indipendentemente dallo stato iniziale.
